\subsection{A second direct arithmetic implementation of negation}
Let $p(a,b)=\frac{(a+b)(a+b+1)}2+b$ be the Cantor pairing function
used by the source program. In the published encoding, the implication
constructor is $I(a,b)=4p(a,b)$ and falsity has code $1$. Therefore
$N(t)=I(t,1)=4p(t,1)$.

The elementary identity
\[
 p(t,1)=\frac{(t+1)(t+2)}2+1
       =\frac{t(t+1)}2+t+2=p(0,t)+2
\]
gives the alternative algorithm
\[
        N(t)=4p(0,t)+8.
\]
This identity changes the arithmetic work performed, rather than merely the
names or textual encoding of instructions.

In the existing physical kernel, \texttt{node64} (native entry state $106$)
executes the pairing and multiplication core without the preceding stack
pop. Its clean-scratch contract, with accumulator $T=t$, auxiliary $X=0$,
and auxiliary $Y=0$, is
\[
 (T,W,X,Y)=(t,w,0,0)
   \longmapsto (4p(0,t),w,0,0).
\]
All other registers are preserved. The existing \texttt{node61}
(native entry state $190$) adds four to $T$. Thus the three macro tokens
\[
 \texttt{node64};\quad \texttt{node61};\quad \texttt{node61}
\]
implement the same accumulator and stack transformation as
\[
 \texttt{push};\quad\texttt{incT};\quad\texttt{cons},
\]
which pushes $t$, loads $1$, pops $t$, and forms $4p(t,1)$.
The new implementation uses $24$ primitive source sites, compared with
$47$ for the original stack-based sequence and $20$ for the earlier
move-to-$X$ negation optimization. Fewer primitive sites alone do not
imply fewer physical control states: sharing with the rest of the decoder
is decisive.

The zero-scratch hypothesis must be retained. For $X=x$ and $Y=0$, the
new pairing core instead produces $4p(x,t)+8$, whereas the earlier direct
negation block gives $4p(t+x,1)$. Their difference is $4x$.
Consequently, the weaker $Y=0$ precondition sufficient for the earlier
negation replacement does not justify this replacement.

The appropriate source-level invariant is available at macro boundaries:
\texttt{M6/Source/AXMacro.lean}'s \texttt{registers} definition sets both
logical registers $6$ and $7$ to zero, and the \texttt{Realizes} contracts
in \texttt{M6/Source/AXMacroSound.lean} preserve that full clean store.
The core contract follows from \texttt{M6/Source/AXPair.lean}, theorem
\texttt{pair}, composed with the multiplication contract in
\texttt{M6/Source/AXTimes4.lean}. Binding the moved blocks to a particular
new source tree and checking its physical table remain necessary.
A zero/positive abstraction of individual primitive steps need not retain
this stronger macro invariant; failure of that abstraction to force
register $6$ to zero is not a proof that the actual macro-boundary invariant
fails.

\paragraph{Preloading the pair accumulator gives an 18-site version.}
The same proof yields a shorter implementation. The more general contract
of \texttt{node64}, obtained by retaining the arbitrary initial accumulator
in \texttt{AXPair.pair\_loop}, is
\[
 (T,W,X,Y)=(t,w,x,y)
 \longmapsto (4(y+p(x,t)),w,0,0).
\]
Thus, from the clean-scratch store, two increments of register $7$ followed
by \texttt{node64} produce $4(2+p(0,t))=4p(t,1)$.
The token sequence
\[
  \operatorname{inc}(7);\quad\operatorname{inc}(7);\quad\texttt{node64}
\]
uses existing native entries $82,82,106$ and has only $18$ primitive source
sites. No extra arithmetic-engine state is required. A mixed version
$\operatorname{inc}(7);\texttt{node64};\texttt{node61}$ has $21$ sites
and the same clean-store effect.

These variants are the three choices $y=0,1,2$ in the formula
\[
 4\bigl(y+p(0,t)\bigr)+4(2-y)=4p(t,1).
\]
For $y=0$ the suffix uses two existing four-increment macros; for $y=1$ it
uses one; for $y=2$ it is empty. Each whole variant is treated as one macro
with clean scratch at its entry and exit. Intermediate nonzero values of
register $7$ inside the macro are intentional.
